\documentclass[10pt,a4paper]{amsart}
\usepackage[margin=19mm]{geometry}
\usepackage{amsmath,amssymb,mathtools}
\usepackage[scr=rsfs,cal=euler]{mathalfa}
\usepackage{xfrac}
\usepackage[unicode,hidelinks,bookmarks=false]{hyperref}
\newcommand{\ie}{i.e.\ }
\newcommand{\cf}{cf.\ }
\newcommand{\resp}{respectively\ }
\newcommand{\Subsection}{\subsection}
\providecommand{\nobreakdash}{\nobreak\hbox{-}}
\setlength{\emergencystretch}{2em}
\multlinegap0pt
\usepackage{bm}
\usepackage{bbm}
\usepackage{dsfont}
\usepackage{xspace}
\usepackage{cancel}
\usepackage{mathtools}
\usepackage{amsthm}
\makeatletter
\@ifundefined{theorem}{\newtheorem{theorem}{Theorem}}{}
\@ifundefined{lemma}{\newtheorem{lemma}[theorem]{Lemma}}{}
\@ifundefined{proposition}{\newtheorem{proposition}[theorem]{Proposition}}{}
\makeatother
\title[Nilpotent primitive linear groups over finite fields]{Nilpotent primitive linear groups over finite fields}
\author{A. S. Detinko, D. L. Flannery}
\date{Source: arXiv:2507.20628v1 (CC BY 4.0)}
\begin{document}
\maketitle

\noindent\emph{Source and licence.} Nilpotent primitive linear groups over finite fields. Version arXiv:2507.20628v1 (CC BY 4.0), distributed under the Creative Commons Attribution 4.0 International licence, as recorded in the arXiv licence metadata for this exact version and shown on the abstract page https://arxiv.org/abs/2507.20628v1. Full rights notice: licenses/arxiv-2507.20628-CC-BY-4.0.txt.\par\smallskip

\noindent\textit{Editorial scope.} Counted unit: the Theorem whose source label is \texttt{firstclassif} in the article (source0.tex lines 831--855, source label \texttt{firstclassif}), delivered together with the article's own complete original proof of it (source0.tex lines 856--1049, one \texttt{proof} environment of 194 lines). No mathematics is written, repaired or re-derived: statement and proof are the article's own text line for line. Required context retained uncounted: normcentr (lines 175--184); sizeprim (lines 216--221); likesim (lines 237--246); likersim (lines 262--267); konzal (lines 296--309); evenimprim (lines 653--659); prea2c (lines 749--762); a2cirred (lines 776--787). Every definition, display and label of those lines is retained unchanged. Omitted from this package and not redistributed: 1--174, 185--215, 222--236, 247--261, 268--295, 310--652, 660--748, 763--775, 788--830, 1050--1310. Nothing inside a retained excerpt is omitted or altered: every retained body is delivered exactly as the article's lines are written, so no omission receipt of any kind accompanies this package. Citations: the retained excerpts carry no citation command at all, so no bibliography entry of the article is transcribed and no reference is redistributed. Reference adaptation: none was needed and none was made; every reference inside the delivered unit resolves inside it, so no unresolved reference or citation is produced. Printed slips kept literal: no printed word, symbol, hypothesis or number of the counted statement or of its proof was altered. Layout and class: the article's own class is \texttt{amsart} with packages including amssymb, euscript; the wrapper is the lane's pinned \texttt{amsart} editorial wrapper at 10pt on A4, which restates the theorem-like environments the retained text needs and adds the article's own macro definitions that the retained text needs (none), with extra wrapper packages (bm, bbm, dsfont, xspace, cancel, mathtools). The delivered numbering is the unit's own. Assets: no retained span contains a figure environment, a graphics-inclusion command, a PDF-page inclusion, a table or a TikZ picture; no image file is copied into this package, the package declares no asset dependency and no image file is redistributed with it. Licence: version arXiv:2507.20628v1 of this article is distributed under the Creative Commons Attribution 4.0 International licence, as recorded in the arXiv licence metadata for this exact version and shown on the abstract page https://arxiv.org/abs/2507.20628v1. A full rights notice is delivered at licenses/arxiv-2507.20628-CC-BY-4.0.txt. Characters: every retained excerpt is pure ASCII, so the delivered engine, XeLaTeX with its default Latin Modern text font, reports no missing character.
\begin{proposition}
\label{normcentr}
 Let $C$ be an irreducible subgroup of $\Delta^{\! \times}$.
 The $\mathrm{GL}(n,\mathbb{F})$-centraliser of 
 $C$ is $\Delta^{\! \times}$, and the 
 $\mathrm{GL}(n,\mathbb{F})$-normaliser of 
 $C$ is the semidirect product of   
 $\Delta^{\! \times}$ with a cyclic group of order
 $n$.
\end{proposition}

\begin{lemma}
\label{sizeprim}
 For each integer $r>1$ dividing $k$, $A$  
 has a system of imprimitivity consisting of $r$
 components.
\end{lemma}

\begin{proposition}
\label{likesim}
 Let $k$ be prime, and denote $\ker \phi$ by $A'$.
\begin{itemize}
\item[{\rm (i)}]  Each $V_i$   
 is a faithful irreducible $A'$-module.
\item[{\rm (ii)}]  
 $|A|$ divides $k(q^{n/k}-1)$.
\end{itemize}
\end{proposition}

\begin{proposition}
\label{likersim}
 A subgroup $C$ of  $\Delta^{\! \times}$  is 
 imprimitive if and only if $|C|$ divides 
 $k(q^{n/k}-1)$ for some prime $k$ dividing  $n$.
\end{proposition}

\begin{theorem}
\label{konzal}
 Suppose $G$ is nilpotent primitive.  
\begin{itemize}
%\item[{\rm (i)}] $[G,G]$ is cyclic,
\item[{\rm (i)}] Every odd order subgroup of $G$ is cyclic.
\item[{\rm (ii)}] If $G_2$ is nonabelian then 
 $q\equiv 3$ mod $4$, and $G_2$ has maximal nilpotency 
 class: either $G_2 \cong Q_8$,
 or $G_2$ has order at least $16$ 
 and is generalised quaternion, 
 dihedral, or semidihedral. 
\end{itemize}
\end{theorem}

\begin{theorem}
\label{evenimprim}
 Let $G= G_2 \times C$ be nonabelian nilpotent primitive, 
 and let $A$ be a cyclic subgroup of $G$ of index $2$.
 If $A$ is imprimitive then every $A$-system of  
 imprimitivity consists of $2$ components. 
\end{theorem}

\begin{lemma}
\label{prea2c}
% Let $\mathbb{E}$ be any field, $B$ an abelian
% irreducible  subgroup of $\mathrm{GL}(r,\mathbb{E})$,
% and $\bar{B}\leq B$. The irreducible $\bar{B}$-submodules 
% of the natural module $W$ for $\mathrm{GL}(r,\mathbb{E})$
% have dimension $|\langle \bar{B} \hspace*{.2mm} 
% \rangle_{\mathbb{E}} : \mathbb{E}1_r|$, 
% and they are all isomorphic to each other.
If $\bar{A}\leq A$ then the irreducible $\bar{A}$-submodules
 of $V$ have dimension $|\langle \bar{A} \hspace*{.2mm} 
 \rangle_{\mathbb{F}} : \mathbb{F}1_n|$, 
and they are all isomorphic to each other.
\end{lemma}

\begin{lemma} 
\label{a2cirred}
\begin{itemize}
\item[{\rm (i)}] 
 An irreducible $A_2$-submodule of $V$ has
 dimension $2$.
\item[{\rm (ii)}] 
 $\mathbb{K} = \langle C\rangle_\mathbb{F}$, and an 
 irreducible $C$-submodule of $V$ has dimension 
 $m$.
\end{itemize}
\end{lemma}

\begin{theorem}
\label{firstclassif}
% Let $G = G_2 \times C \leq  \mathrm{GL}(2m, \mathbb{F})$, 
% where $C =\langle c\rangle$ is odd order cyclic, 
% $G_2 = \langle a, g\rangle$ is the Sylow 
% $2$-subgroup of $G$, $A_2 = \langle a\rangle$ is a maximal 
% abelian subgroup of $G_2$ of index $2$, 
% $A = A_2 \times C$
%% = \langle c, a_2\rangle$ 
% is irreducible, and  
% $|\langle C\rangle_\mathbb{F} : \mathbb{F}1_n| = m$. 
% Assume the notation and hypotheses introduced before 
% Lemma$~\ref{}$.
 $G$ but not $A$ is primitive if and only if $A$ has an 
 imprimitivity system 
%$\mathcal{S} = \{ V_1 ,  V_2\}$,
 of size $2$,  
 $C$ is conjugate to $\langle (x,x) \rangle$ 
%$(C_1, C_1)$ 
 for some primitive subgroup 
%$C_1= 
$\langle x \rangle$ of 
 $\mathrm{GL}(m, \mathbb{F})$, and 
 $a^2 = g^2 = [a, g] = -1_n$.
\end{theorem}

\begin{proof}
% We note at the outset that if $K$ is any subgroup of
% $A$ then $\langle K \rangle_{\mathbb{F}}$ is a field.
 Suppose $G$ but not $A$ is primitive.
 By Theorem~\ref{evenimprim},
 $A$ has an imprimitivity system $\mathcal{S} = \{ V_1, V_2 \}$.
% By irreducibility of $A$, $\mathrm{dim}\,  V_1$ $=$
% $\mathrm{dim}\,  V_2=m$.
% Recall the notation 
 Let $\phi$ be the permutation representation of $A$ in 
 $\mathrm{Sym}(2)$ arising from the action on $\mathcal{S}$, 
 and set $\ker \phi=A'$.
 Since $|C|$ is odd, necessarily $C \leq A'$, 
and then by
%  $C$ $=$
 %$\mathrm{diag} (C_1, C_2)$, $C_i \leq \mathrm{GL}(m, \mathbb{F})$, 
% $\{ \mathrm{diag} (\rho_1(c), \rho_2(c)) \mid c \in C \}$, 
% $C_i \leq \mathrm{GL}(m, \mathbb{F})$, 
% up to conjugacy. By 
 Lemmas~\ref{prea2c} and \ref{a2cirred},
%, $V_1$ and
% $V_2$ are irreducible $C$-modules, and because
% (clearly) $V_1 \cong V_2$ as $C$-modules, in fact 
% $C_1$ is $\mathrm{GL}(m, \mathbb{F})$-conjugate to $C_2$. 
% Hence 
% $C_1 = C_2$ is irreducible, and we may take 
% $x=y$. 
 $C = \langle (x,x) \rangle$ up to conjugacy,
 where $C_1 = \langle x \rangle \leq \mathrm{GL}(m, \mathbb{F})$ 
 is irreducible. Since $\langle C\rangle_\mathbb{F}$ $\leq$
 $\langle A'\rangle_\mathbb{F}$
 and $|\Delta : \langle C\rangle_\mathbb{F}| =2$, 
%$$
%2m = |\Delta: \mathbb{F}1_n| = 
%  |\Delta : \langle C\rangle_\mathbb{F}|
%  |\langle C\rangle_\mathbb{F} : \mathbb{F}1_n|  
%$$ 
 either $\langle A' \rangle_\mathbb{F}$ $=$ 
 $\Delta$ or $\langle A'\rangle_\mathbb{F} = 
 \langle C\rangle_\mathbb{F}$. 
 If $\langle A' \rangle_\mathbb{F} = 
 \Delta$ then $A'$ is irreducible, 
 which is definitely not true.
 Thus $\langle A'\rangle_\mathbb{F} = 
 \langle C\rangle_\mathbb{F}$.    
 An element $h\neq 1$ of $A' \cap A_2$ has 2-power order;
 but $h$ $\in$  $\langle C\rangle_\mathbb{F}$, which
 does not contain an element of order 4, so $h = -1_n$. 
 %$|h| = 2$. 
 %%We have that $h \in \langle (x, y)\rangle$
 %Write $h = (x,x)$ for some 
 %$x \in \langle C_1 \rangle_{\mathbb{F}}$, $|x| = 2$. 
 %%$\langle C_1 \rangle_{\mathbb{F}}$ contains all scalars
 %%of $\mathrm{GL}(m,\mathbb{F})$, so the element of 
 %%order 2 in  $\langle C_1 \rangle_{\mathbb{F}}$ is 
 %In fact $x =  -1_m$ because $C_1$ 
 %is irreducible.
 %singer cycle contains all scalars
 %Consequently $h \in \langle -1_n\rangle$, i.e.
 %$A_2 \cap A' \leq \langle -1_n\rangle$. 
 Consequently 
 $A' = (A' \cap A_2)\times C = \langle -1_n, C \rangle$ and 
 $A'|_{V_i} = \langle -1_m, C_1 \rangle$. Since 
 $A'|_{V_i}$ is primitive (otherwise we get an $A$-system of
 imprimitivity of size greater than 2, violating 
 Theorem~\ref{evenimprim}), $C_1$ is primitive.
 From $|A: A'|=2$ and $|A'| = 2|C|$ we infer $|A_2| =4$
 and $|G_2|=8$. 
 By Theorem~\ref{konzal}, then, $G_2 \cong Q_8$.
 %($D_8$ has a noncyclic abelian normal subgroup, 
 %which is normal in $G$). Note that $-1_n \in G_2$.
 We have seen that $|a|=4$, and 
 $g \notin Z(G_2) = \langle -1_n \rangle$
 implies 
%$|g| = 4$. Hence 
 $a^2=g^2= [a,g] = -1_n$.

 %Now suppose $A$ has imprimitivity system $\mathcal{S}$,
 %$C$ and $G_2 \cong Q_8$ are as described and 
 Now we prove the converse. Suppose 
 %$A$ has imprimitivity system $\mathcal{S}$, 
 $C = \langle (x,x) \rangle$ where 
 % (C_1, C_1)$ for a 
 $\langle x \rangle := C_1$ is a primitive subgroup of 
 $\mathrm{GL}(m, \mathbb{F})$, and $G_2 \cong Q_8$.
% and $-1_n \in G_2$. 
 Also suppose $\mathcal{T}$ $=$ $\{  W_1 , \ldots \, \! ,  W_k\}$
 is a $G$-system of imprimitivity, and let $\varphi
 : G \rightarrow \mathrm{Sym}(k)$ be the associated 
%transitive 
 permutation representation.
 Of course, $\mathcal{T}$ is simultaneously an $A$-system of 
 imprimitivity, so that
% for every subgroup of $G$, and in particular 
 if $k$ is divisible by an odd prime $t$
 then by Lemma~\ref{sizeprim}, $A$ has an imprimitivity 
 system of size $t$.
 By Proposition~\ref{likesim}, $|A| = 4|C_1|$ divides 
 $t(q^{m/t}-1)(q^{m/t}+1)$.  
 Some prime divisor $r$ of $|C_1|$ must divide $q^{m/t}+1$
 by Proposition~\ref{likersim}, and therefore $r$ divides $q^m+1$.
 %As $C_1$ $\leq$ $\mathrm{GL}(m,\mathbb{F})$ is irreducible,
 However $|C_1|$ divides $q^m-1$, so $r$ divides $2q^m$, 
 a contradiction. 
%It follows that $k$ is a power of 2. 
% Since $k$ divides $2m$, $k=2$.
 Hence $k=2$.
% By Lemma~\ref{no7}, 
% 
%As before it can be shown that 
 %$\langle A \cap \ker \varphi \rangle_{\mathbb{F}}
 %= \langle C \rangle_{\mathbb{F}}$, which has 
 %no element of  
%
%is the $\mathbb{F}$-enveloping 
% algebra of $A$ $\cap$ $\ker \varphi$. 
%
% order 4, and thus $A_2 \cap \ker \varphi$ $=$
% $\langle -1_n\rangle$. 
%
% (which is a subgroup of $G_2$ by hypothesis). 
 
 Since
% $\varphi(C)$ is trivial, 
 $\varphi(A_2)$ is transitive, after replacing $G$ 
 with a suitable conjugate we have $a = (u, v) w$
 where $u, v$ $\in$ $\mathrm{GL}(m, \mathbb{F})$ and
$$
 w = 
 \left( \begin{array}{cc}
 0_m & 1_m \\
  1_m & 0_m 
 \end{array}
 \right) ,
$$
 %and $C\leq A'$ is $2\times 2$ block diagonal
 %form, with the $m\times m$ blocks conjugate 
 %in $\mathrm{GL}(m,\mathbb{F})$. A further conjugation leaves
 %$a$ in the same form, 
 with $C$ retaining its original form
 %$\mathrm{diag}(C_1,C_1)$. 
 $\langle (x,x) \rangle$. The relation 
 $a^2 = -1_n$ yields $v$ $=$ $-u ^{-1}$.
% $a = (u, -u^{-1})T_2$. 

 Recall Proposition~\ref{normcentr}.
 Since $a$ centralises $C$, $u$ centralises $C_1$,
 and therefore $u \in \langle C_1\rangle_\mathbb{F}$.
% $\mathrm{GL}(m,\mathbb{F})$-centraliser of the 
% irreducible subgroup $C_1$ of  $\mathrm{GL}(m,\mathbb{F})$,
% viz. $\langle C_1\rangle_\mathbb{F}$. Since
% As $g$ normalises $A$, it normalises 
% $\Delta^{\! \times}$.
%= \langle C, a\rangle_\mathbb{F}$.
 Note that $d =  (1_m, -1_m)$ inverts $a$, so
  it normalises but does not centralise 
 $\Delta^{\! \times}$. Certainly $g$ normalises  
 $\Delta^{\! \times}$, so
%$e \not \in \Delta$, 
  $g = db$ for some $b \in \Delta^{\! \times}$.
% by Proposition~\ref{normcentr}. 
%Since $\mathcal{T}$ is a $G$-system of imprimitivity, either
 
 Over $\langle C \rangle_\mathbb{F}$, $\Delta$ has basis
 $\{ 1_n, a\}$.
 Thus if 
%$h \in \Delta$ acts trivially on $\mathcal{T}$ then 
%$h \in \langle C \rangle_\mathbb{F}$.
% 
% Either $g \in \ker \varphi$ or $|\varphi (g)| = 2$.
 $g \in \ker \varphi$ then $b$ $=$ 
 $(y, y)$ for some
 %$b$ $=$ 
 %$(v, v) \in \langle C\rangle_\mathbb{F}$, 
 %i.e. 
 $y\in \langle C_1\rangle_\mathbb{F}$. Since $g^2 = -1_n$
 it follows that $y^2 = -1_m$,
% That is, $|y| =4$, 
 which is impossible.

 Suppose $g \notin \ker \varphi$, so that
% i.e. $\langle g\rangle$ acts on $\mathcal{T}$ transitively. 
 %$e$ fixes each of $W_1$ and $W_2$, 
 %$b \in \Delta$ must interchange them, i.e. 
 $b$ $=$ $(u y ,-u^{-1} y)w$ for some 
 $y \in \langle C_1\rangle_\mathbb{F}$.
% so $g$ $=$ $(u x ,u^{-1} x)w$.
 %Then $g$ $=$ $eb$ $=$ 
 %$(v u , v u^{-1})T_2$. Since 
 Once again, $g^2 = -1_n$ implies 
 $\langle C_1\rangle_\mathbb{F}$ has an element
% (here, $u$)???
 of order 4. This final contradiction proves that $G$ is primitive. 
\end{proof}

\end{document}
